\section{22 September 2026: Livingston outcome vintage and Wallabout allocation}

\paragraph{Historical outcomes need a dated source rule.}
The historical route selects cohorts by initial filing date but reads proposed
Class~A units and lot identifiers from Housing Database release 25Q4. The
official 23Q4 archive records 370 Livingston's job B00781447 at 144 units,
compared with 99 in 25Q4. Its initial filing date is 28 September 2022.
The recorded declaration 2026010400010003, physical page~16, reproduces an
April 2025 HPD application with 99 units and intended 485-x Option~B treatment.
The current 99 is therefore not an observed pre-policy count. These sources
do not identify the exact amendment date or establish a causal policy effect.

The 23Q4 file also records 372 Livingston's B00814498 at 105 units, filed
120 days after 370. Both contemporaneous records share lot~16 and explicitly
refer to the neighboring building under a separate application. They support
one earlier two-building proposal with 249 units in that release. Current
historical linking sees later lots~22 and~16, and treats the common-lot recovery
as a later transaction; it misses the link. The January 2025 zoning description
2025021001245002, page~2, draws the subsequent three-building development on
old lots~16/26 and tentative lots~16/22/26. It does not establish the original
proposal's land allocation or earlier floor allocation.

\paragraph{Diagnostic scope.}
The new audit holds current panel selection and membership fixed and compares
unique DOB jobs with the active and inactive-inclusive 23Q4 files. Of 1,904
historical constituents, 1,902 match and 225 have different unit counts,
affecting 218 current parents. In the current 554-parent weighting sample,
591 of 592 constituents match; 109 differ, affecting 106 parents. Four of the
six current historical exact-99 constituents were not 99 in 23Q4:
Livingston 144, Lafayette 104, East 198th Street 64, and Vyse Avenue 54.
The other two were 99 in both releases. The active-file-only check finds 220
changes; the inactive-inclusive file adds five further changes and 24 matches.

These are vintage differences, which may reflect amendments or administrative
corrections. They are not estimates of policy responses. Reconstructing the
pre-policy distribution requires a common outcome vintage, contemporaneous
link fields, and explicit treatment of inactive and unmatched records. This
follow-up saves the diagnostic but does not change production units or links.
The historical counterfactual needs that correction before further estimation.

\paragraph{Wallabout implementation.}
The completed map overlay allocates the three cohort buildings' recorded
39,323 square feet across common prefiling 18v2\_1 parcels: 2,507.69 on old
lot~37, 18,342.52 on old~41, and 18,472.79 on old~122. Each successor's
recorded area is distributed by its GIS overlap shares. This is an approximate
zoning allocation; the ground total is administrative. Old lot~23 contributes
no ground. The later lot~37 belongs to separate 251 Wallabout and has a
different boundary from the earlier lot with that number.

All three old parcels record zero building floor. The production manual table
now applies these rows: parent ground falls from 68,805 to 39,323 and
residential FAR rises from about 4.70 to 5.03. Only this parent's covariates
change. Constituent units, dates, membership and refiling fields are identical
to the prior build. The number of flagged earlier-floor estimates remains six.
The boundary screen now has 111 unresolved and 753 supported parents;
Wallabout is the only changed pass/flag status. The weighting samples remain
554 historical and 310 post-policy parents. The land correction alone moves
the weighted historical exact-99 share from 1.3537427\% to 1.3536772\%, holding
the current outcome-vintage rule fixed. The calibration moment error is below
$6.5\times10^{-10}$; 493 of 499 bootstrap draws succeed, with six calibration
failures recorded. These refreshed outputs remain provisional historical
comparisons pending the source-vintage correction above.

\paragraph{Other evidence and remaining work.}
The actual 70-page Rockaway condominium plans, CRFN 2022000210978, are now
saved. Pages~57--58 locate Building~D's units 18/19/20 and its limited common
elements. Their floor areas are not a measure of total development ground;
outdoor/common-ground allocation and Phase~IV still require review. The
310/322 Concourse investigation fixes 29,609.6 square feet of land, while its
dated drawings leave an earlier-building allocation to reconcile. Broadway's
unfinished earlier structure and Charles's earlier/newer proposal relationship
also remain evidence tasks, rather than numerical choices ready for approval.

\paragraph{Sources and reproduction.}
The source manifest is in the acquisition task's
\path{code/site_research_next10_2026-09-22/queued_followup_sources.csv}.
Root \texttt{make dof-site-review} acquires the checksum-pinned DCP archive
and runs \path{compare_hdb_unit_vintages.R} and
\path{measure_wallabout_allocation.R}. They produce five audit tables with
SaveData reports. Wallabout's production dependency is the committed manual
parcel table; no main task depends on an audit. The detailed evidence and
implementation review is \path{report/queued_followup_2026-09-22.md} in the
parcel audit task.
